\clearpage
\VerdierPhysicalPageAnchor{0190}%
\begingroup\small\setlength{\parskip}{0pt}\setlength{\emergencystretch}{1.5em}
\setlength{\abovedisplayskip}{4pt plus 2pt minus 2pt}%
\setlength{\belowdisplayskip}{4pt plus 2pt minus 2pt}%

\VerdierSemanticAnchor{ch3d-2-1-5}
\noindent\textbf{注 2.1.5.}\quad
仍在引理 (2.1.4) 的假设下，并进一步假设
$X,Y\in\operatorname{Ob}\operatorname{comp}^{-}(\mathcal A)$。上面的论证表明，
必要时可将 $X$ 替换为一个与之同伦等价的复形，从而假定 $m$ 是满态射。上述 $X_n$ 列的
构造于是表明，归纳系 $X_n$ 是定常的（即 $X_n=Y$ 且
$u_n=\operatorname{id}$ 对足够小的 $n$ 成立），并且对每个 $n$，
$X_n\in\operatorname{Ob}\operatorname{comp}^{-}(\mathcal A)$。

\VerdierSemanticAnchor{ch3d-2-2}
\subsection*{2.2. 某些对象集的研究。}

\VerdierSemanticAnchor{ch3d-2-2-1}
\noindent\textbf{2.2.1.}\quad
设 $M$ 是 $\mathcal A$ 的一个对象集。我们考虑下列性质：

\VerdierSemanticAnchor{ch3d-2-2-1-e1}
\noindent $(E_1)$ 零对象属于 $M$。两个属于 $M$ 的对象之直和仍属于 $M$。

\VerdierSemanticAnchor{ch3d-2-2-1-e2}
\noindent $(E_2)$ 对每个满态射 $X\xrightarrow{u}L$，其中 $X$ 是
$\mathcal A$ 的对象且 $L$ 是 $M$ 的元素，存在满态射
$v\colon L'\longrightarrow L$，其中 $L'\in M$，并且该态射通过 $u$ 分解。

\VerdierSemanticAnchor{ch3d-2-2-1-e3}
\noindent $(E_3)$ 对每个整数 $n$ 及每条正合列
\[
0\longrightarrow X\longrightarrow L_n\longrightarrow L_{n-1}
\longrightarrow\cdots\longrightarrow L_0\longrightarrow0,
\qquad L_i\in M,\quad 0\le i\le n,
\]
存在满态射 $L'\longrightarrow X$，其中 $L'\in M$。

\VerdierSemanticAnchor{ch3d-2-2-1-e4}
\noindent $(E_4)$ 对每个满态射 $u\colon L\longrightarrow L'$，其中
$L,L'\in M$，$u$ 的核属于 $M$。

\VerdierSemanticAnchor{ch3d-2-2-1-e5}
\noindent $(E_5)$ 对每个扩张
$0\longrightarrow L\longrightarrow Z\longrightarrow L'\longrightarrow0$，其中
$L,L'\in M$，对象 $Z$ 属于 $M$；$0$ 属于 $M$。

\VerdierSemanticAnchor{ch3d-2-2-1-e6}
\noindent $(E_6)$ 集合 $M$ 是正合的，即每当 $\mathcal A$ 的一条短正合列中有
两个对象属于 $M$ 时，第三个对象也如此；$0$ 属于 $M$。

\VerdierSemanticAnchor{ch3d-2-2-1-implications}
注意有蕴含 $(E_6)\Rightarrow(E_4)$ 与 $(E_5)$、
$(E_4)\Rightarrow(E_3)$ 以及 $(E_5)\Rightarrow(E_1)$。一条无界正合列

\VerdierSemanticAnchor{ch3d-2-2-1-infinite-presentation}
\[
\cdots\longrightarrow L_n\longrightarrow L_{n-1}\longrightarrow\cdots
\longrightarrow L_0\longrightarrow X\longrightarrow0,
\qquad L_i\in M,\quad i\in\mathbb N,
\]
称为一个\emph{$\infty$-$M$-表示}，其所表示的对象为 $X$。设 $M$ 是满足
性质 $(E_1)$、$(E_2)$ 与 $(E_3)$ 的对象集。以 $M^{\wedge}$ 表示如下对象集：
对象 $X$ 属于 $M^{\wedge}$，如果它具有一个 $\infty$-$M$-表示，并且对于
每个整数 $n$ 和每条正合列

\VerdierSemanticAnchor{ch3d-2-2-1-m-wedge}
\[
0\longrightarrow X'\longrightarrow L_n\longrightarrow L_{n-1}
\longrightarrow\cdots\longrightarrow L_0\longrightarrow X\longrightarrow0,
\qquad L_i\in M,\quad 0\le i\le n,
\]
对象 $X'$ 都有一个 $\infty$-$M$-表示。
\endgroup
